\documentclass[../../../../main.tex]{subfiles}

\begin{document}

\subsection*{第1問}

[解] $n$ 個から $2k$ 個をとり出し, これを $2$ つずつの組に分ける場合の数と同じである。

したがって, $n \ge 2k$ のもとで
\[
\frac{{}_n\mathrm{C}_{2k} \cdot {}_{2k}\mathrm{C}_2 \cdot {}_{2k-2}\mathrm{C}_2 \cdots {}_4\mathrm{C}_2 \cdot {}_2\mathrm{C}_2}{k!}
\]
\[
= \frac{\frac{n!}{2k!(n-2k)!} \cdot \frac{2k(2k-1)}{2} \frac{4\cdot 3}{2} \cdots \frac{2\cdot 1}{2}}{k!}
\]
\[
= \frac{n!}{k!(n-2k)!} \cdot \frac{1}{2^k} \quad \text{通り}
\]
であり, $n < 2k$ の時は $0$ 通り。以上まとめて,
\[
f(n, k) = 
\begin{cases}
\frac{n!}{(n-2k)! \cdot k!} \frac{1}{2^k} \, \text{通り} & (n \ge 2k) \\
0 \, \text{通り} & (n < 2k)
\end{cases}
\]

次に, $f(n, k) = f(n, 1)$ について, $n \ge 2$ から $f(n, 1) \neq 0$ だから, $n \ge 2k$ のもとで考えれば良い。
\[
f(n, k) = f(n, 1)
\]
\[
\iff \frac{1}{2^k} \frac{n!}{k!(n-2k)!} = \frac{1}{2} \frac{n!}{(n-2)!}
\]
\[
\iff 2^{k-1} k! (n-2k)! = (n-2)!
\]
\[
\iff 2^{k-1} = \frac{(n-2)(n-3) \dots (n+1-2k)}{k!} \quad \dots \text{①}
\]
ここで, 連続した整数は $k!$ で割り切れるから, $A \in \mathbb{N}$ として,
\[
(n-2) \cdots (n+1-2k) = A(2k-2)!
\]
とおける。①から
\[
\text{①} \iff 2^{k-1} = A \cdot (2k-2)(2k-3) \cdots (k+1) \quad \dots \text{②}
\]
同様に, $(2k-2) \dots (k+1)$ は $(k-2)!$ で割り切れるから, $k \ge 5$ の時 $2^{k-1}$ が $3$ で割り切れることになり不適。又, $k=2, 4$ の時は ②から, $2^{k-1}$ が各々 $3$ ($k=2$), $5$ ($k=4$) で割り切れることになり不適。以上から $k=3$ が必要。この時①から,
\[
3 \cdot 2 \cdot 2^2 = (n-2) \cdots (n-5)
\]
\[
1 \cdot 2 \cdot 3 \cdot 4 = (n-2) \cdots (n-5)
\]
となり, $n=6$ とすれば十分。以上から, もとめるのは $(n, k) = (6, 3)$ である。


\subsection*{第2問}

[解] 正六角形の一辺の長さは $1$ であり, 三角形ができる時, その面積は
\begin{enumerate}
  \item[$1^\circ$] $A_1 A_2 A_6$ 形 $\dots \frac{\sqrt{3}}{4}$
  \item[$2^\circ$] $A_1 A_2 A_5$ 形 $\dots \frac{\sqrt{3}}{2}$
  \item[$3^\circ$] $A_1 A_3 A_5$ 形 $\dots \frac{3}{4}\sqrt{3}$
\end{enumerate}
のいずれかである。$3$ つの頂点のえらび方 $6^3$ 通りが同様にたしからしい。以下, $1$ つ目に $A_1$ がえらばれたとして

$1^\circ$ の時
\begin{quote}
どの点が頂点になるかで $2 \times 3 = 6$ 通り
\end{quote}

$2^\circ$ の時
\begin{quote}
$1$ つの頂点に対し, のこりの頂点が $2$ 通りあるから
\[
4 \times 2 + 2 \times 2 = 12 \text{ 通り}
\]
\end{quote}

$3^\circ$ の時
\begin{quote}
$2$ 通り
\end{quote}

だから, これらの $6$ 倍ずつえらび方があるので, もとめる期待値として
\[
E = \frac{6}{36} \cdot \frac{\sqrt{3}}{4} + \frac{12}{36} \cdot \frac{\sqrt{3}}{2} + \frac{2}{36} \cdot \frac{3}{4}\sqrt{3}
\]
\[
= \frac{\sqrt{3}}{36} \left( \frac{3}{2} + 6 + \frac{3}{2} \right)
\]
\[
= \frac{\sqrt{3}}{4}
\]

\begin{center}
\begin{tikzpicture}[scale=1.5]
  % Hexagon vertices
  \coordinate (A1) at (90:1.8);
  \coordinate (A2) at (30:1.8);
  \coordinate (A3) at (-30:1.8);
  \coordinate (A4) at (-90:1.8);
  \coordinate (A5) at (-150:1.8);
  \coordinate (A6) at (150:1.8);

  \draw (A1) -- (A2) -- (A3) -- (A4) -- (A5) -- (A6) -- cycle;

  % Internal lines from A1
  \draw[dashed] (A1) -- (A3);
  \draw[dashed] (A1) -- (A4);
  \draw[dashed] (A1) -- (A5);

  % Vertex labels
  \node[above] at (A1) {$A_1$};
  \node[above right] at (A2) {$A_2$};
  \node[right] at (A3) {$A_3$};
  \node[below] at (A4) {$A_4$};
  \node[left] at (A5) {$A_5$};
\end{tikzpicture}
\end{center}


\subsection*{第3問}

[解] $Q(t, t^2)$ における法線ベクトルが $\vec{n} = \begin{pmatrix} -2t \\ 1 \end{pmatrix}$ だから

法線上の点 $P(X, Y)$ は $s \in \mathbb{R}$ として
\[
\begin{pmatrix} X \\ Y \end{pmatrix} = \begin{pmatrix} t \\ t^2 \end{pmatrix} + s \begin{pmatrix} -2t \\ 1 \end{pmatrix} \quad \dots \text{①}
\]
とかける。$\left| \begin{pmatrix} -2t \\ 1 \end{pmatrix} \right| = \sqrt{4t^2+1}$ から, $\overline{PQ} = \sqrt{4t^2+1} \, |s|$ である。

(ロ)から $s > 0$ だから, (イ)に代入して
\[
\sqrt{4t^2+1} \, s = \sqrt{4t^2+1} (t - t^2)
\]
$\sqrt{4t^2+1} > 0$ から $s = t - t^2 \quad \dots \text{②}$ なので, ②を①に代入して $P$ は
\[
\begin{cases}
X = 2t^3 - 2t^2 + t \\
Y = t
\end{cases} \quad (0 \le t \le 1)
\]
で表されるパラメータ曲線の上うち, $Y \ge 0$ の部分をうごく。$t$ をけして
\[
X = 2Y^3 - 2Y^2 + Y \quad (0 \le Y \le 1) \equiv f(Y)
\]
このグラフの試形は右図で, もとめるのは斜線部の面積で,
\[
S = \int_0^1 (\sqrt{Y} - f(Y)) \, dY
\]
\[
= \left[ \frac{2}{3} Y^{\frac{3}{2}} - \frac{1}{2} Y^4 + \frac{2}{3} Y^3 - \frac{1}{2} Y^2 \right]_0^1
\]
\[
= \frac{1}{3}
\]

$\left[ x = 2y^3 - 2y^2 + y \text{ は 3次関数なので, そのグラフの面積はばらさずに計算してよいだろう} \right]$

\begin{center}
\begin{tikzpicture}[scale=2.2]
  % Axes
  \draw[->] (-0.2,0) -- (1.3,0) node[right] {$x$};
  \draw[->] (0,-0.2) -- (0,1.3) node[above] {$y$};
  \node[below left] at (0,0) {$0$};

  % Functions: y = x^2 (x = sqrt(y)) and x = f(y) = 2y^3 - 2y^2 + y
  \draw[domain=0:1,samples=50] plot ({2*\x^3 - 2*\x^2 + \x}, {\x}) node[below right] {$f(y)$};
  \draw[domain=0:1,samples=50] plot ({\x}, {\x*\x}) node[above left] {$y=x^2$};

  % Fill shaded area
  \fill[pattern=north east lines,pattern color=gray!60] 
    plot[domain=0:1,samples=50] ({2*\x^3 - 2*\x^2 + \x}, {\x}) -- 
    plot[domain=1:0,samples=50] ({\x}, {\x*\x}) -- cycle;

  % Key points
  \draw[dashed] (1,0) node[below] {$1$} -- (1,1) -- (0,1) node[left] {$1$};
  \fill (1,1) circle (0.6pt);
\end{tikzpicture}
\end{center}


\subsection*{［角度の扱いメモ］}

\begin{enumerate}
  \item[(1)] $\cos\theta \dots 0 \le \theta \le \frac{\pi}{2}$ が明らかな時, 有効
  \item[(2)] $\tan\theta \dots$ これで立式して, $\tan \leftrightarrow \cos$ のりかえを行う (推奨)
  \item[(3)] $\tan\theta = \frac{S}{C} = \frac{\triangle ABC}{\vec{a} \cdot \vec{b}}$ (推奨)
\end{enumerate}

［$\tan\theta \leftrightarrow \cos\theta$ のりかえ］
$\frac{1}{\cos^2}, \frac{1}{\sin^2}$ がでてきたら速攻のりかえる。


\subsection*{第4問}

[解] $0 \le \alpha < \pi/2 \quad \dots \text{①}$

$Z > 0 \quad \dots \text{②}$ から, $P$ は $A, B, C$ であることに適する。
\[
\vec{AP} = \begin{pmatrix} x-1 \\ y-1 \\ Z \end{pmatrix}, \quad 
\vec{BP} = \begin{pmatrix} x-1 \\ y+1 \\ Z \end{pmatrix}, \quad 
\vec{CP} = \begin{pmatrix} x \\ y \\ Z \end{pmatrix}
\]
(ロ)～(ニ)より, $\vec{n} = \begin{pmatrix} 0 \\ 0 \\ 1 \end{pmatrix}$ として,
\[
\begin{cases}
\vec{n} \cdot \vec{AP} = |\vec{n}| |\vec{AP}| \cos\pi/4 \\
\vec{n} \cdot \vec{BP} = |\vec{n}| |\vec{BP}| \cos\pi/4 \\
\vec{n} \cdot \vec{CP} = |\vec{n}| |\vec{CP}| \cos\alpha
\end{cases}
\iff
\begin{cases}
Z = \sqrt{(x-1)^2 + (y-1)^2 + Z^2} \cdot \frac{1}{\sqrt{2}} \\
Z = \sqrt{(x-1)^2 + (y+1)^2 + Z^2} \cdot \frac{1}{\sqrt{2}} \\
Z = \sqrt{x^2 + y^2 + Z^2} \cos\alpha
\end{cases}
\]
①②から, 両辺 2 乗しても同値。
\[
\begin{cases}
Z^2 = (x-1)^2 + (y+1)^2 = (x-1)^2 + (y-1)^2 & \dots \text{②} \\
Z^2 \cdot \tan^2\alpha = x^2 + y^2 & \dots \text{③}
\end{cases}
\]
②の右側の不等式から, $(y+1)^2 = (y-1)^2 \iff y=0 \quad \dots \text{④}$ だから, ②③に代入して
\[
\begin{cases}
Z^2 = (x-1)^2 + 1 = x^2 - 2x + 2 & \dots \text{⑤} \\
Z^2 \cdot \tan^2\alpha = x^2 & \dots \text{⑥}
\end{cases}
\]
したがって, ①～② $\land$ ④ $\land$ ⑤ $\land$ ⑥ をみたす $P$ の数をもとめれば良い。⑤から,

$1^\circ$ $\alpha \neq 0$ の時

⑥から $Z^2 = \left( \frac{x}{\tan\alpha} \right)^2$ だから, ⑤に代入して, $t = \tan\alpha$ とおくと,
\[
\begin{cases}
x \neq 0 \\
(t^2-1) x^2 - 2t^2 x + 2t^2 = 0 \quad \dots \text{⑦}
\end{cases}
\]
⑦をみたす $x$ に対し, ⑥から $Z$ が $Z = \left|\frac{x}{\tan\alpha}\right|$ と一意に定まり, $y=0$ だから, ⑦をみたす $x$ の数が $P$ の数に等しい。$t^2 - 1 \neq 0$ の時, ⑦の判別式 $D$ として
\[
D/4 = t^4 - (t^2-1) \cdot 2t^2 = t^2 (2 - t^2)
\]
だから,
\[
\begin{cases}
D/4 > 0 \iff 2 > t^2 \text{ の時 } 2 \text{ 個} \\
D/4 = 0 \iff 2 = t^2 \text{ の時 } 1 \text{ 個} & \dots A_1 \\
D/4 < 0 \iff 2 < t^2 \text{ の時 } 0 \text{ 個}
\end{cases}
\]
である ($\alpha \neq 0$ から, $x=0$ は解にならない)。

一方, $t^2 - 1 = 0 \iff \alpha = \pi/4$ ($\because$ ①) の時, ⑦ $\to x = 1$, ⑥から $Z=1$ となって $1$ 個。

$2^\circ$ $\alpha = 0$ の時

⑥から $x=0$ だから ⑤より $Z = \sqrt{2}$ となる。

以上まとめて, ⑦の解が ($D \ge 0$ の時) $x = \frac{t^2 \pm \sqrt{t^2(2-t^2)}}{t^2 - 1}$ であるから
\[
Z = \left| \frac{x}{t} \right| = \left| \frac{t \pm \sqrt{2-t^2}}{t^2 - 1} \right|
\]
であることをあわせて,
\[
\begin{cases}
\alpha = 0 \text{ の時}, Z = \sqrt{2} \text{ の } 1 \text{ つ} \\
0 < \tan\alpha < \sqrt{2} \land \alpha \neq \pi/4 \text{ の時}, Z = \left| \frac{t \pm \sqrt{2-t^2}}{t^2 - 1} \right| \text{ の } 2 \text{ つ} \\
\tan\alpha = \sqrt{2} \text{ の時}, Z = \sqrt{2} \text{ の } 1 \text{ つ} \\
\sqrt{2} < \tan\alpha \text{ の時}, 0 \text{ 個} \\
\alpha = \pi/4 \text{ の時}, Z = 1 \text{ の } 1 \text{ つ}
\end{cases}
\]

\begin{center}
\begin{tikzpicture}[scale=1.5, >=stealth]
  % Axes
  \draw[->] (0,0,0) -- (2.5,0,0) node[right] {$y$};
  \draw[->] (0,0,0) -- (0,2.5,0) node[above] {$z$};
  \draw[->] (0,0,0) -- (0,0,2.5) node[below left] {$x$};

  % Origin / Point C
  \node[below left] at (0,0,0) {$C$};

  % Points A, B, P
  \coordinate (C) at (0,0,0);
  \coordinate (A) at (1,0,1);
  \coordinate (B) at (-1,0,1);
  \coordinate (P) at (0,1.8,0.8);

  % Draw points and lines
  \draw[dashed] (C) -- (P) node[midway, right] {$\alpha$};
  \draw[dashed] (A) -- (P);
  \draw[dashed] (B) -- (P);
  \draw (0,0,0) -- (1,0,1) node[right] {$A$};
  \draw (0,0,0) -- (-1,0,1) node[left] {$B$};

  \fill (P) circle (1.5pt) node[above] {$P$};
  \fill (A) circle (1.5pt);
  \fill (B) circle (1.5pt);
  \fill (C) circle (1.5pt);
\end{tikzpicture}
\end{center}


\subsection*{第5問}

[解] (1) 図のような展開図での頂点を $A_1, A_2, \dots, A_n$ とする。
又, 底面の各点を $B_1 \dots B_n$ とする。$A_k$ から $B_k B_{k+1}$ ($k=1, 2 \dots n, B_{n+1} = B_1$) に下ろした垂足を $H_k$ とする。
$OH_1 = x, A_1 H_1 = y$ とおく ($0 < x < y$)。

$B_k B_{k+1} = 2x \tan \frac{\pi}{n}$

さて, 展開図が円に内接するとき
\[
A_1 O = -x + y = 1 \iff y = 1 - x \quad \left(0 < x < \frac{1}{2}\right)
\]
である。立体の高さとして
\[
h = \sqrt{y^2 - x^2} = \sqrt{1 - 2x}
\]
又, 底面積 $S$ は
\[
S = n \times (\triangle O B_1 B_2) = n \times \frac{1}{2} \cdot x \cdot \left(2x \tan\frac{\pi}{n}\right) = n x^2 \tan\frac{\pi}{n}
\]
だから体積 $V$ は
\[
V = \frac{1}{3} S \cdot h = \frac{1}{3} n \tan\frac{\pi}{n} \, x^2 \sqrt{1 - 2x}
\]
$x$ の部分を $f(x)$ とおくと, $f(x) \ge 0$ から $f(x)^2$ が最大の時 $f(x)$ も最大。
\[
(f(x))^2 = 4x^3 - 10x^4 = 2x^3 (2 - 5x)
\]
より下表を考える。

\[
\begin{array}{c|c|c|c|c|c}
x & 0 & \dots & \frac{2}{5} & \dots & \frac{1}{2} \\ \hline
f' & & + & 0 & - & \\ \hline
f^2 & & \nearrow & \text{極大} & \searrow &
\end{array}
\]

よって $x = 2/5$ の時 $f(x)$ も最大で
\[
V = V_n = \frac{1}{3} n \tan\frac{\pi}{n} \cdot \left(\frac{2}{5}\right)^2 \sqrt{\frac{1}{5}} = \frac{4\sqrt{5}}{375} n \tan\frac{\pi}{n}
\]

(2) $t = \frac{\pi}{n}$ とおくと $n \to \infty$ で $t \to 0$ で
\[
V_n = \frac{4\sqrt{5}}{375} \pi \frac{\tan t}{t} \longrightarrow \frac{4\sqrt{5}}{375} \pi \quad (n \to \infty, t \to 0)
\]

\begin{center}
\begin{tikzpicture}[scale=1.3]
  \coordinate (O) at (0,0);
  \coordinate (A1) at (90:2);
  \coordinate (A2) at (30:2);
  \coordinate (A3) at (-30:2);

  \coordinate (B1) at (120:0.8);
  \coordinate (B2) at (60:0.8);
  \coordinate (B3) at (0:0.8);

  \draw (O) circle (2);

  \draw (A1) -- (B1) -- (B2) -- cycle;
  \draw (A2) -- (B2) -- (B3) -- cycle;
  \draw (O) -- (B1);
  \draw (O) -- (B2);

  \coordinate (H1) at ($(B1)!0.5!(B2)$);
  \draw[dashed] (O) -- (H1) node[midway, left] {$x$};
  \draw[dashed] (A1) -- (H1);

  \node[above] at (A1) {$A_1$};
  \node[above right] at (A2) {$A_2$};
  \node[right] at (A3) {$A_3$};
  \node[left] at (B1) {$B_1$};
  \node[above] at (B2) {$B_2$};
  \node[right] at (B3) {$B_3$};
  \node[below left] at (O) {$O$};
  \node[above right] at (H1) {$H_1$};
\end{tikzpicture}
\qquad
\begin{tikzpicture}[scale=1.3]
  \draw (0,0) node[below left] {$O$} -- (1.5,0) node[below right] {$H_1$} node[midway, below] {$x$} -- (1.5,2.5) node[above] {頂点} -- cycle node[midway, above left] {$y$};
  \node[right] at (1.5,1.25) {$h$};
  \draw (1.3,0) -- (1.3,0.2) -- (1.5,0.2);
\end{tikzpicture}
\end{center}


\subsection*{第6問}

[解] (1) (ロ)から, $f(x) = a(x+1)(x-1)(x-\alpha)$ ($\alpha \in \mathbb{R}$) とおける。($a \neq 0$ の時)

$1^\circ$ $0 \le x$ の時

(イ) $\iff f(x) \ge 1 - x \iff a(x+1)(x-\alpha)(x-1) \ge -(x-1)$ で, $x=1$ で等号成立, $0 \le x < 1$ の時は $x-1 < 0$ だから
\[
a(x+1)(x-\alpha) \le -1
\]
\[
a x^2 + a(1-\alpha) x + 1 - a\alpha \le 0 \quad \dots \text{①}
\]
この左辺を $g(x)$ とおく。

$\text{i)}$ $a > 0$ の時

$0 \le x \le 1$ で $g(x) \le 0$ であれば良い。条件は
\[
g(0) \le 0 \land g(1) \le 0 \iff 1 - a\alpha \le 0 \land 2a - 2a\alpha + 1 \le 0
\]

$\text{ii)}$ $a < 0$ の時

$g(x)$ の軸 $x = \frac{\alpha - 1}{2}$ で場合分け, $g(x) = 0$ の判別式を $D$ とおく。
\[
\begin{cases}
0 \le \frac{\alpha-1}{2} \le 1 \text{ の時}, D \le 0 \iff a^2(1-\alpha)^2 - 4a(1-a\alpha) \le 0 \iff a \left[ a(\alpha+1)^2 - 4 \right] \le 0 \dots \star_1 \\
\text{otherwise} \quad g(0) \le 0 \land g(1) \le 0
\end{cases}
\]

$2^\circ$ $x \le 0$ の時

$f(x) \ge 1+x \iff a(x-1)(x-\alpha)(x+1) \ge (x+1)$ で, $1^\circ$ と同様に
\[
a(x-1)(x-\alpha) \ge 1
\]
\[
a x^2 - a(1+\alpha) x + a\alpha - 1 \ge 0 \quad \dots \text{②}
\]
この左辺を $h(x)$ とおく。

$\text{i)}$ $a > 0$ の時

軸 $x = \frac{\alpha+1}{2}$ で場合分け, $h(x)=0$ の判別式を $D'$ とおく。
\[
\begin{cases}
-1 \le \frac{\alpha+1}{2} \le 0 \text{ の時}, D' \le 0 \iff a^2(\alpha+1)^2 - 4a(a\alpha-1) \ge 0 \iff a \left[ a(\alpha-1)^2 + 4 \right] \ge 0 \dots \star_2 \\
\text{otherwise} \quad g(-1) \ge 0 \land g(0) \ge 0 \iff 2a + 2a\alpha - 1 \ge 0 \land a\alpha - 1 \ge 0
\end{cases}
\]

$\text{ii)}$ $a < 0$ の時

条件は $g(-1) \ge 0 \land g(0) \ge 0$

ここで, $a < 0$ から $\star_1 \iff a(\alpha+1)^2 - 4 \ge 0$ (これをみたす $\alpha, a$ はない)。
同じく $a > 0$ から $\star_2 \iff a(\alpha-1)^2 + 4 \ge 0 \iff \alpha, a$ は任意。

だから, $1^\circ, 2^\circ$ をまとめて,

$a > 0$ の時
\[
\begin{cases}
1 - a\alpha \le 0 \land 2a - 2a\alpha + 1 \le 0 \land (\alpha \le -3 \text{ or } 0 \le \alpha) \land \\
a\alpha - 1 \ge 0 \land 2a + 2a\alpha - 1 \ge 0
\end{cases}
\]
$a < 0$ の時
\[
\begin{cases}
a\alpha - 1 \ge 0 \land 2a + 2a\alpha - 1 \ge 0 \land (\alpha \le 1 \text{ or } 3 \le \alpha) \land \\
1 - a\alpha \le 0 \land 2a - 2a\alpha + 1 \le 0
\end{cases}
\]

ここで, $f(x) = a(x^2-1)(x-\alpha) = a(x^3 - \alpha x^2 - x + \alpha)$ から
\[
b = -a\alpha, \quad d = a\alpha, \quad c = -a \quad \dots \text{③}
\]
である。$\alpha$ は, $a>0$ の時 $1 < \alpha$, $a<0$ の時 $\alpha < -1$ なので条件はまとめて
\[
\begin{cases}
\alpha = \frac{b}{-a} + 1 \\
\alpha = \frac{d}{a}
\end{cases} \quad \text{で} \quad
\begin{cases}
1 - a\alpha \le 0 \land 2a - 2a\alpha + 1 \le 0 \\
2a + 2a\alpha - 1 \ge 0
\end{cases} \quad \dots \text{④}
\]
一方, $a=0$ の時, $f(x)$ が $1$ 次以下で不適なので $b \neq 0$ のもとで $f(x) = b(x^2-1)$ とおける。
(ハ)の条件は $f(0) \ge 1 \iff 1+b \le 0$ である。

以上まとめて,
\[
b = -d, \quad c = -a, \quad 1 + b \le 0, \quad
\begin{cases}
2a + 2b + 1 \le 0 \\
2a - 2b - 1 \ge 0
\end{cases}
\]

(2) (1)から, $f(x) = ax^3 + bx^2 - ax - d$ とおける。$f'(x) = 3ax^2 + 2bx - a$ から
\[
\{f'(x) - x\}^2 = (3ax^2)^2 + (2b-1)^2 x^2 + a^2 + 2 \left[ 3a(2b-1) x^3 - 3a^2 x^2 - (2b-1)ax \right]
\]
\[
= 9a^2 x^4 + \left[ (2b-1)^2 - 6a^2 \right] x^2 + a^2 + (\text{奇数次項})
\]
より,
\[
\int_{-1}^1 \{f'(x) - x\}^2 \, dx = 2 \int_0^1 \left( 9a^2 x^4 + \left[ (2b-1)^2 - 6a^2 \right] x^2 + a^2 \right) dx
\]
\[
= 2 \left[ a^2 \int_0^1 (9x^4 - 6x^2 + 1) dx + (2b-1)^2 \int_0^1 x^2 dx \right]
\]
\[
= 2 \left[ \frac{4}{5} a^2 + \frac{1}{3}(2b-1)^2 \right] \equiv A \quad \dots \text{⑤}
\]
ここで, $(a, b)$ は領域をうごくから,
\[
\begin{cases}
0 \le a \le \frac{1}{2} \text{ の時}, b \le -\frac{1}{2} \\
\frac{1}{2} \le a \text{ の時}, b \le -a - \frac{1}{2}
\end{cases}
\]
したがって
\[
|a| \le \frac{1}{2} \text{ の時}, A \text{ は } (a, b) = (0, -1) \text{ で } \min 6
\]
\[
\frac{1}{2} \le |a| \text{ の時}, A \text{ は } 0 \le a \le \frac{1}{2} \text{ の時より大きい}
\]
から, $A$ は $(a, b) = (0, -1)$ で $\min 6$ をとる。この時
\[
f(x) = -x^2 + 1
\]

[別] 軽〜く処理するなら…
(ハ)の条件は $f'(1) \le -1, f'(-1) \ge 1, f(0) \ge 1$ である。

\begin{center}
\begin{tikzpicture}[scale=1.5]
  % Axes for (a,b) plane
  \draw[->] (-1.5,0) -- (1.5,0) node[right] {$a$};
  \draw[->] (0,-2) -- (0,1) node[above] {$b$};
  \node[above left] at (0,0) {$0$};

  % Line b = -1/2 and b = -a - 1/2
  \draw[dashed] (-1.5,-0.5) -- (1.5,-0.5) node[right] {$b = -\frac{1}{2}$};
  \draw[dashed] (0.5,-0.5) -- (1.5,-1.5) node[right] {$b = -a - \frac{1}{2}$};
  \draw[dashed] (-0.5,-0.5) -- (-1.5,-1.5);

  % Shaded region
  \fill[pattern=north east lines, pattern color=gray!60] (-0.5,-2) rectangle (0.5,-0.5);
  \fill[pattern=north east lines, pattern color=gray!60] (0.5,-0.5) -- (1.5,-1.5) -- (1.5,-2) -- (0.5,-2) -- cycle;
  \fill[pattern=north east lines, pattern color=gray!60] (-0.5,-0.5) -- (-1.5,-1.5) -- (-1.5,-2) -- (-0.5,-2) -- cycle;

  % Key points
  \draw[dashed] (0.5,0) node[above] {$\frac{1}{2}$} -- (0.5,-0.5);
  \draw[dashed] (-0.5,0) node[above] {$-\frac{1}{2}$} -- (-0.5,-0.5);
  \fill (0,-1) circle (1.5pt) node[left] {$(0, -1)$};
\end{tikzpicture}
\end{center}

\end{document}
